\PassOptionsToPackage{table,xcdraw}{xcolor}
\documentclass[10pt, a4paper]{article}

\usepackage[T1]{fontenc}
\usepackage[utf8]{inputenc}
\usepackage{tgpagella}
\usepackage{tgcursor}
\usepackage[spanish, es-noshorthands]{babel}
\usepackage{amsmath, amssymb}
\usepackage{siunitx}
\usepackage{booktabs}
\usepackage{tabularx}
\usepackage{xcolor}
\usepackage{tikz}
\usetikzlibrary{shapes.geometric, positioning}
\usepackage[american]{circuitikz}
\usepackage[most]{tcolorbox}
\usepackage[margin=1.5cm]{geometry}
\usepackage{fancyhdr}
\usepackage{lastpage}

\setlength{\headheight}{16pt}
\setlength{\emergencystretch}{2em}

\definecolor{ucbblue}{RGB}{0, 51, 102}
\definecolor{accentteal}{RGB}{0, 102, 102}
\definecolor{neutralbg}{RGB}{246, 248, 251}
\definecolor{warnbg}{RGB}{254, 246, 246}
\definecolor{warnframe}{RGB}{153, 30, 30}

\pagestyle{fancy}
\fancyhf{}
\lhead{\textbf{IMT-221 --- Circuitos Electrónicos II}}
\chead{Formulario de Ecuaciones -- Examen EC2}
\rhead{Hasta Lab 6}
\cfoot{\footnotesize Página \thepage\ de \pageref{LastPage} \quad|\quad Hoja de fórmulas sin ejemplos resueltos (uso sujeto a autorización docente)}
\renewcommand{\headrulewidth}{0.5pt}
\renewcommand{\footrulewidth}{0.4pt}

\newtcolorbox{formulabox}[2][ucbblue]{%
    colback=white,
    colframe=#1,
    colbacktitle=#1,
    coltitle=white,
    fonttitle=\bfseries\small,
    title=#2,
    boxrule=0.7pt,
    arc=1.5mm,
    left=7pt,
    right=7pt,
    top=5pt,
    bottom=5pt,
    before skip=6pt,
    after skip=6pt
}

\begin{document}

% =========================================================================
% ENCABEZADO
% =========================================================================
\begin{center}
    {\large\textbf{\textcolor{ucbblue}{IMT-221 --- Formulario de Ecuaciones para el Examen EC2}}}\\[0.08cm]
    {\footnotesize\textbf{Alcance:} amplificadores operacionales reales, In-Amp/CMRR, diodos y rectificación de precisión, filtros activos de segundo orden (hasta Lab 6).}
\end{center}

\begin{tcolorbox}[colback=neutralbg, colframe=ucbblue!70, boxrule=0.6pt, arc=1.2mm, left=6pt, right=6pt, top=4pt, bottom=4pt]
    \footnotesize
    \textbf{Uso:} Hoja de fórmulas de referencia, \textbf{sin ejemplos resueltos ni resultados de los problemas}. \textbf{Su autorización para la prueba corresponde al docente:} \quad $\square$ Autorizado \quad $\square$ No autorizado \qquad \textbf{V.\textsuperscript{o} B.\textsuperscript{o}:} \rule{3.0cm}{0.4pt}
\end{tcolorbox}

% =========================================================================
% SECCION 1: AMPLIFICADORES OPERACIONALES REALES
% =========================================================================
\begin{formulabox}{1. Amplificadores operacionales reales}
    \begin{minipage}[t]{0.62\textwidth}
        \textbf{Amplificador no inversor} \textit{(realimentación negativa; operación lineal):}
        \begin{align}
            A_{CL} &= \frac{V_o}{V_i} = 1 + \frac{R_f}{R_g}, \label{eq:acl} \\[2pt]
            f_{BW} &\approx \frac{f_T}{|A_{CL}|} \qquad \text{(polo dominante; } \text{noise gain} = A_{CL}\text{)}, \label{eq:fbw}
        \end{align}
        donde $f_T$ es la frecuencia de ganancia unitaria o $\mathrm{GBW}$ en el modelo indicado.

        \vspace{0.18cm}
        \textbf{Respuesta sinusoidal y \textit{slew rate} (SR):}
        \begin{align}
            v_o(t) &= V_p \sin(2\pi f t + \varphi), \label{eq:seno} \\[2pt]
            \mathrm{SR}_{\mathrm{req}} &= \max\left|\frac{dv_o}{dt}\right| = 2\pi f V_p, \qquad V_{p,\max} \approx \frac{\mathrm{SR}}{2\pi f}. \label{eq:sr}
        \end{align}
    \end{minipage}%
    \hfill
    \begin{minipage}[t]{0.35\textwidth}
        \centering
        \vspace{0pt}
        \begin{circuitikz}[>=stealth, thick, scale=0.72, transform shape]
            \draw (0,0) node[op amp, yscale=-1] (op) {};
            \draw (op.+) to[short, -o] ++(-1.1,0) node[left]{$v_i$};
            \draw (op.out) to[short, *-o] ++(1.0,0) node[right]{$v_o$};
            \draw (op.-) -- ++(0,-0.85) coordinate(ninv);
            \draw (ninv) to[R, l_=$R_g$] ++(0,-1.45) node[ground]{};
            \draw (ninv) to[short, *-] ++(0.3,0) to[R, l=$R_f$] ++(1.7,0) -| (op.out);
        \end{circuitikz}\\[0.1cm]
        {\scriptsize\textit{Esquema simbólico no inversor (ref. $0\,\mathrm{V}$)}}
    \end{minipage}

    \vspace{0.15cm}
    \hrule
    \vspace{0.12cm}
    \footnotesize
    \textbf{Conversiones:} $1\,\mathrm{V/\mu s} = 10^6\,\mathrm{V/s}$; \quad $V_{pp} = 2 V_p$ (para una senoide sin \textit{offset}).\\
    \textbf{Condición de verificación:} Comprobar \textbf{separadamente} banda ($f_{BW}$), $\mathrm{SR}$, rieles/excursión permitida y rango de modo común. Los límites específicos se suministran en cada enunciado.
\end{formulabox}

% =========================================================================
% SECCION 2: ACONDICIONAMIENTO DIFERENCIAL E INSTRUMENTACION
% =========================================================================
\begin{formulabox}{2. Acondicionamiento diferencial e instrumentación}
    \begin{minipage}[t]{0.62\textwidth}
        \begin{align}
            v_d &= v_+ - v_-, \qquad v_{cm} = \frac{v_+ + v_-}{2}, \label{eq:vd_vcm} \\[3pt]
            v_o &= A_d v_d + A_{cm} v_{cm} + V_{\mathrm{offset}} \qquad \text{(modelo lineal)}, \label{eq:vo_inamp} \\[3pt]
            \mathrm{CMRR} &= \left|\frac{A_d}{A_{cm}}\right|, \qquad \mathrm{CMRR}_{\mathrm{dB}} = 20\log_{10}(\mathrm{CMRR}), \label{eq:cmrr} \\[3pt]
            \Delta v_o &= A_d \Delta v_d + A_{cm} \Delta v_{cm} \qquad \text{(\textit{offset} y ganancias constantes)}. \label{eq:delta_vo}
        \end{align}
    \end{minipage}%
    \hfill
    \begin{minipage}[t]{0.35\textwidth}
        \centering
        \vspace{0.25cm}
        \begin{tikzpicture}[>=stealth, thick, scale=0.82, transform shape]
            \draw[fill=neutralbg, draw=ucbblue, rounded corners=2pt] (0,0) rectangle (2.6,1.8);
            \node[align=center, font=\footnotesize\bfseries, text=ucbblue] at (1.3, 1.15) {In-Amp /\\Diferencial};
            \node[align=center, font=\scriptsize] at (1.3, 0.4) {$A_d,\; A_{cm}$};
            \draw[->] (-0.9, 1.35) node[left]{$v_+$} -- (0, 1.35) node[right, font=\scriptsize]{$+$};
            \draw[->] (-0.9, 0.45) node[left]{$v_-$} -- (0, 0.45) node[right, font=\scriptsize]{$-$};
            \draw[->] (2.6, 0.9) -- (3.5, 0.9) node[right]{$v_o$};
            \draw (1.3, 0) -- (1.3, -0.45) node[ground]{};
        \end{tikzpicture}\\[0.05cm]
        {\scriptsize\textit{Bloque lineal simbólico (ref. común)}}
    \end{minipage}

    \vspace{0.12cm}
    \hrule
    \vspace{0.12cm}
    \footnotesize
    \textbf{Unidades y caso ideal:} $A_d$ y $A_{cm}$ son ganancias de tensión en $\mathrm{V/V}$. Si $A_{cm} = 0$, el $\mathrm{CMRR}$ ideal se considera infinito.
\end{formulabox}

% =========================================================================
% SECCION 3: DIODOS Y RECTIFICACION
% =========================================================================
\begin{formulabox}{3. Diodos y rectificación}
    \begin{itemize}
        \setlength{\itemsep}{2pt}
        \item \textbf{Diodo ideal:} conduce en sentido directo y bloquea en inverso (modelo por tramos).
        \item \textbf{Rectificador convencional} de media onda positiva con carga a $0\,\mathrm{V}$ \textit{(modelo de umbral)} y \textbf{referencia ideal}:
        \begin{equation}
            v_o(t) \approx \max\big[v_i(t) - V_\gamma,\; 0\big], \qquad v_{\mathrm{ref}}(t) = \max\big[v_i(t),\; 0\big]. \label{eq:rect}
        \end{equation}
        \item \textbf{Superdiodo no inversor idealizado:} con diodo conductor y realimentación negativa cerrada, $v_o \approx v_i$; si el diodo bloquea, la realimentación se abre y \textbf{no se puede imponer} $v_+ \approx v_-$.
        \item \textbf{Error absoluto medio de $N$ muestras} \textit{(solo si se solicita):}
        \begin{equation}
            E_{\mathrm{MAE}} = \frac{1}{N}\sum_{k=1}^{N} \big|v_o[k] - v_{\mathrm{ref}}[k]\big|. \label{eq:mae}
        \end{equation}
    \end{itemize}

    \begin{center}
        \begin{circuitikz}[>=stealth, thick, scale=0.76, transform shape]
            % Convencional
            \node[font=\footnotesize\bfseries, text=ucbblue] at (1.5, 2.35) {Convencional (media onda $+$)};
            \draw (0,0) node[ground]{} to[sV] (0,1.65) -- (0.55,1.65) to[D, l=$D$] (2.35,1.65) to[short, *-o] (3.0,1.65) node[right]{$v_o$};
            \node[left, font=\footnotesize] at (-0.45, 0.85) {$v_i$};
            \node[below, font=\tiny] at (1.05, 1.50) {A};
            \node[below, font=\tiny] at (1.85, 1.50) {K};
            \draw (2.35,1.65) to[R, l=$R_L$] (2.35,0) node[ground]{};

            % Superdiodo no inversor (op.- arriba, op.+ abajo -> sin cruces)
            \node[font=\footnotesize\bfseries, text=ucbblue] at (9.4, 2.35) {Superdiodo no inversor (idealizado)};
            \draw (7.4, 0.85) node[op amp] (op) {};
            \draw (op.+) to[short, -o] ++(-1.0,0) node[left]{$v_i$};
            \draw (op.out) to[short, *-] ++(0.55,0) coordinate(va) node[below]{\scriptsize $v_a$} to[D, l=$D$] ++(2.0,0) coordinate(vo) to[short, *-o] ++(0.75,0) node[right]{$v_o$};
            \node[below, font=\tiny] at ([xshift=0.6cm, yshift=-0.12cm]va) {A};
            \node[below, font=\tiny] at ([xshift=1.4cm, yshift=-0.12cm]va) {K};
            \draw (vo) to[R, l=$R_L$] (vo |- 0,0) node[ground]{};
            \draw (op.-) -- ++(0,0.75) -| (vo);
        \end{circuitikz}
    \end{center}

    \vspace{0.05cm}
    \hrule
    \vspace{0.10cm}
    \footnotesize
    \textbf{Alcance del parámetro $V_\gamma$:} $V_\gamma$ depende del modelo especificado; \textbf{no} se considera una constante universal del diodo real.
\end{formulabox}

\newpage

% =========================================================================
% SECCION 4: FILTROS ACTIVOS DE SEGUNDO ORDEN
% =========================================================================
\begin{formulabox}{4. Filtros activos de segundo orden}
    \begin{minipage}[t]{0.66\textwidth}
        \textbf{Función de transferencia canónica pasa-bajos:}
        \begin{equation}
            H(s) = \frac{V_o(s)}{V_i(s)} = \frac{K\omega_0^2}{s^2 + (\omega_0/Q)s + \omega_0^2}. \label{eq:hs_canon}
        \end{equation}
        \textbf{Frecuencia característica y normalizada:}
        \begin{equation}
            \omega_0 = 2\pi f_0, \qquad f_0 = \frac{\omega_0}{2\pi}, \qquad x = \frac{\omega}{\omega_0} = \frac{f}{f_0}, \label{eq:w0_x}
        \end{equation}
        con $\omega$ en $\mathrm{rad/s}$ y $f$ en $\mathrm{Hz}$.
    \end{minipage}%
    \hfill
    \begin{minipage}[t]{0.31\textwidth}
        \centering
        \vspace{0.35cm}
        \begin{tikzpicture}[>=stealth, thick, scale=0.85, transform shape]
            \draw[->] (-0.9, 0) node[left]{$V_i(s)$} -- (0, 0);
            \draw[fill=neutralbg, draw=ucbblue, rounded corners=2pt] (0, -0.65) rectangle (2.6, 0.65);
            \node[align=center, font=\small\bfseries, text=ucbblue] at (1.3, 0.12) {$H(s)$};
            \node[align=center, font=\scriptsize] at (1.3, -0.35) {$K,\; \omega_0,\; Q$};
            \draw[->] (2.6, 0) -- (3.5, 0) node[right]{$V_o(s)$};
        \end{tikzpicture}\\[0.1cm]
        {\scriptsize\textit{Bloque lineal (sin carga mutua)}}
    \end{minipage}

    \vspace{0.15cm}
    \textbf{Respuesta en frecuencia} \textit{(sustituyendo $s = j\omega$, con $j^2 = -1$):}
    \begin{align}
        |H(j\omega)| &= \frac{|K|}{\sqrt{(1 - x^2)^2 + (x/Q)^2}}, \label{eq:mag} \\[3pt]
        G_{\mathrm{dB}} &= 20\log_{10}\left|\frac{V_o}{V_i}\right| = 20\log_{10}|H(j\omega)|, \label{eq:gdb} \\[3pt]
        \varphi(\omega) &= \arg\big[H(j\omega)\big] = -\operatorname{atan2}\!\left(\frac{x}{Q},\; 1 - x^2\right) \qquad \text{(para $K > 0$; selección de cuadrante correcta)}. \label{eq:fase}
    \end{align}

    \textbf{Propiedades en la frecuencia natural, Butterworth y cascada:}
    \begin{itemize}
        \setlength{\itemsep}{2pt}
        \item En la frecuencia natural ($\omega = \omega_0$, $x = 1$): \quad $|H(j\omega_0)| = |K|\,Q$; \quad para $K > 0$, $\varphi(\omega_0) = -90^\circ$.
        \item \textbf{Butterworth pasa-bajos de segundo orden:} \quad $Q = \dfrac{1}{\sqrt{2}}$.
        \item \textbf{Advertencia sobre $-3\,\mathrm{dB}$:} \textbf{No presuponer} que $f_0$ es la frecuencia de $-3\,\mathrm{dB}$; hallar el punto respecto de la ganancia en banda pasante cuando corresponda.
        \item \textbf{Etapas en cascada} \textit{(desacopladas o sin carga significativa):}
        \begin{equation}
            H_{\mathrm{total}}(s) \approx H_1(s)\,H_2(s), \qquad G_{\mathrm{total,dB}}(f) \approx G_{1,\mathrm{dB}}(f) + G_{2,\mathrm{dB}}(f). \label{eq:cascada}
        \end{equation}
    \end{itemize}
\end{formulabox}

% =========================================================================
% SECCION 5: MATEMATICA Y UNIDADES DE USO TRANSVERSAL
% =========================================================================
\begin{formulabox}[accentteal]{5. Matemática y unidades de uso transversal}
    \begin{minipage}[t]{0.52\textwidth}
        \textbf{Relaciones matemáticas:}
        \begin{align}
            \omega &= 2\pi f, \\[2pt]
            \log_{10}(ab) &= \log_{10}(a) + \log_{10}(b) \qquad (a, b > 0), \\[2pt]
            |a + jb| &= \sqrt{a^2 + b^2}, \\[2pt]
            \arg(a + jb) &= \operatorname{atan2}(b, a).
        \end{align}
    \end{minipage}%
    \hfill
    \begin{minipage}[t]{0.44\textwidth}
        \textbf{Prefijos y unidades:}
        \vspace{0.15cm}
        \begin{center}
            \renewcommand{\arraystretch}{1.25}
            \begin{tabularx}{\textwidth}{@{}l X@{}}
                \toprule
                \textbf{Unidad} & \textbf{Equivalencia} \\
                \midrule
                $1\,\mathrm{k\Omega}$ & $10^{3}\,\Omega$ \\
                $1\,\mathrm{M\Omega}$ & $10^{6}\,\Omega$ \\
                $1\,\mathrm{\mu F}$ & $10^{-6}\,\mathrm{F}$ \\
                $1\,\mathrm{mV}$ & $10^{-3}\,\mathrm{V}$ \\
                \bottomrule
            \end{tabularx}
        \end{center}
    \end{minipage}

    \vspace{0.12cm}
    \hrule
    \vspace{0.12cm}
    \footnotesize
    \textbf{Interpretación de decibelios:} Los resultados en $\mathrm{dB}$ corresponden a \textbf{razones de amplitud}, no a una diferencia de voltajes.
\end{formulabox}

% =========================================================================
% NOTA FINAL DE HIPOTESIS Y VALIDEZ
% =========================================================================
\begin{tcolorbox}[colback=warnbg, colframe=warnframe, boxrule=0.7pt, arc=1.2mm, left=7pt, right=7pt, top=5pt, bottom=5pt]
    \footnotesize
    \textbf{Nota general de validez:} Las fórmulas aproximadas requieren que se cumplan sus hipótesis de aplicación. Los valores reales de un circuito integrado (CI) y las condiciones de operación deben provenir exclusivamente del enunciado o del circuito especificado.
\end{tcolorbox}

\end{document}
